% Folder 84, batch 09 — archivists' pages 161–180.
% Pass: Opus 5.5 (claude-opus-5-5), 2026-10-09 — first pass, unchecked against the pages by a human.
%
% The even pages 162–180 are blank versos and are absent. Notation fixed for
% this batch: \mathcal{A} for his script A (algebras A_1, A_2, A), \underline{O}
% for his underlined O (the base ring), \otimes_0 for his ⊗ with subscript 0,
% \otimes^{\circ} for his ⊗ with superscript circle, \varpi for his barred ϖ.
\documentclass[11pt,a4paper]{article}
\input{../preamble/grothendieck}

\folder{84}
\batch{9}
\pages{161}{180}
\foldertitle{[Formes quadratiques] : notes manuscrites (s.d.).}
\dating{[à partir de 1982-vers 1986]}
\watermark{Édition de démonstration}

\begin{document}

\section{Produit tensoriel d'extensions munies de $\chi$-trivialisations}

\page{161}
\note{le calcul qui ouvre la page suppose une suite exacte $0 \to L \to E \to L' \to 0$ et un $\pi$ introduits avant ce lot.}

\[
x \longmapsto \pi'(x) = \pi(x) - (\chi - 1)\,x
\qquad\qquad 0 \to L \to E \to L' \to 0
\]

\struck{\ill{}} $E \to E$ \struck{$\pi'(L)$}
\[
\left\{
\begin{array}{l}
\pi'|L = \text{\struck{\ill{}}}\ \mathrm{id}_L \\
\pi'|E/L = L' \;\;)= -(\chi - 1)\,\mathrm{id}_{L'}
\end{array}
\right.
\iff \pi|L = \chi\,\mathrm{id}_L
\]
\[
\iff \pi : E \to L
\]

\struck{$\pi'(x)$}

\[
x_1 \otimes x_2 + \underbrace{\pi_1'(x_1)}_{u_1} \otimes \pi_2'(x_2)
= u_1 \otimes \bigl(\underbrace{x_2 + \pi_2'(x_2)}_{\chi_2 + \pi_2(x_2) - (\chi - 1) x_2}\bigr)
\]
\[
= (2 - \chi)\,x_2 + \pi_2(x_2) \qquad ?
\]
\note{le « $\chi_2$ » sous l'accolade est lu tel quel ; on attendrait $x_2$. La ligne $(2-\chi)x_2 + \pi_2(x_2)$ est soulignée et suivie d'un point d'interrogation.}

$L_1 \subset E_1$, $L_2 \subset E_2$, avec $\pi_1$ et $\pi_2$ (flèches courbes $E_i \to L_i$) \uncertain{$\chi$-trivialisations} $\Longrightarrow$

\begin{tikzcd}[column sep=small, row sep=small]
E_1 \otimes E_2 \arrow[rr, "?"] \arrow[dr] \arrow[d] & & E_1 \otimes E_2 \arrow[d] \\
E_1 \otimes L_2 \oplus_{L_1 \otimes L_2} L_1 \otimes E_2 \arrow[rr, "{\pi_1, \pi_2}"'] \arrow[drr] & E \arrow[ur] & L_1' \otimes L_2' \\
 & & \overline{L_1' \otimes L_2'}
\end{tikzcd}

\note{lecture du premier diagramme approximative : sous la somme amalgamée on lit « $L_1 \otimes L_2$ » comme indice, et le dernier sommet porte une barre ; la flèche horizontale du haut est marquée « ? ».}

\begin{tikzcd}[column sep=small, row sep=small]
0 \arrow[r] & E_1 \otimes_0 E_2 \arrow[r] \arrow[d, "\pi_2"] & E_1 \otimes E_2 \arrow[r] \arrow[d] & L_1' \otimes L_2' \arrow[r] \arrow[d, no head, "\|" description] & 0 \\
0 \arrow[r] & E_1 \otimes L_2 \arrow[r] \arrow[d, "\text{inc. can.}"] & E \arrow[r] \arrow[d, "?"] & L_1' \otimes L_2' \arrow[r] \arrow[d, "\chi\,\mathrm{id}"] & 0 \\
0 \arrow[r] & E_1 \otimes_0 E_2 \arrow[r] & E_1 \otimes E_2 \arrow[r] & L_1' \otimes L_2' \arrow[r] & {}
\end{tikzcd}

\note{dans la ligne du milieu, le $E_1$ de $E_1 \otimes L_2$ est écrit en surcharge (sur un $L_1$ ?) ; la flèche de droite entre les deux premières lignes est une égalité.}

\struck{$x_1 \quad \pi_1 \quad \alpha +$}

$a_1, \pi_1(x_1) \qquad a_2, \pi_2(x_2)$
\[
\alpha\, x_1 \otimes x_2 + \beta\, x_1 \otimes \pi_2(x_2) + \gamma\, \pi_1(x_1) \otimes x_2 + \delta\, \pi_1(x_1) \otimes \pi_2(x_2)
\]

\page{163}
\[
\pi_i' :\quad
\begin{array}{c|cc}
 & L_i & M_i \\ \hline
L_i & \chi & \tau_i \\
M_i & 0 & 0
\end{array}
\qquad\qquad
\varpi_i' :\quad
\begin{array}{c|cc}
 & L_i & M_i \\ \hline
L_i & 0 & -\tau_i \\
M_i & 0 & \chi
\end{array}
\]

$l_1 \otimes m_2$ \quad \struck{$l_1 + m_2$}

\struck{$\pi_1' \otimes \mathrm{id}_{E_2} + \mathrm{id}_{E_1} \otimes$}

\[
\begin{array}{c|ccc}
1 & 1 & 1 & 1 \\ \hline
\pi_1' & \chi & \chi & \tau_1 \otimes \mathrm{id}_{L_2} \\
\pi_2' & \chi & \mathrm{id}_{L_1} \otimes \tau_2 & \chi \\
\pi_1' \otimes \pi_2' & \chi^2 & \chi\, \mathrm{id}_{L_1} \otimes \tau_2 & \chi\, \tau_1 \otimes \mathrm{id}_{L_2}
\end{array}
\]
\note{les trois colonnes ne portent pas d'en-tête ; la première ligne porte trois « 1 », le premier écrit en surcharge. Les colonnes correspondent vraisemblablement à $L_1 \otimes L_2$, $L_1 \otimes M_2$, $M_1 \otimes L_2$.}

\[
\alpha + \beta_1 \pi_1' + \beta_2 \pi_2' + \gamma\, \pi_1' \otimes \pi_2'
\]
\struck{$\alpha \pi_1' + \beta \pi'$}

\[
\boxed{\alpha + (\beta_1 + \beta_2)\chi + \gamma \chi^2 = \chi}
\]
\[
\alpha + \beta_1 \chi + \underbrace{\beta_2\, \mathrm{id} \otimes \tau_2 + \gamma \chi\, \mathrm{id} \otimes \tau_2}_{(\beta_2 + \gamma\chi)\, \mathrm{id} \otimes \tau_2} = \mathrm{id}_{L_1} \otimes \tau_2
\]
i.e. $\mathrm{id}_{L_1} \otimes \bigl(\beta_1 \chi\, \mathrm{id}_{M_2} + (\beta_2 + \gamma\chi)\tau_2\bigr) = \mathrm{id} \otimes \tau_2$
\note{au-dessus de $\beta_1$, une surcharge « $\alpha+$ » ; lecture incertaine.}

\[
\boxed{\begin{array}{l}
\alpha + \beta_1 \chi = 0 \\
\beta_2 + \gamma \chi = 1
\end{array}}
\]
\note{l'indice de $\beta_1$ dans la première équation encadrée est douteux.}

\page{165}
Pour définir un \uncertain{homomorphisme} $F$, il faut définir

\begin{tikzcd}[column sep=tiny, row sep=small, nodes={font=\scriptsize}]
 & \Sigma = (E_1 \otimes L_2) \oplus (L_1 \otimes E_2) \arrow[dl, "{\mathrm{id}_1 \otimes \alpha_2,\ \alpha_1 \otimes \mathrm{id}_2}"'] \arrow[dr, "{\pi_1 \otimes \mathrm{id}_2,\ \mathrm{id}_1 \otimes \pi_2}"] & \\
E_1 \otimes E_2 \arrow[d, "F_{11}"'] \arrow[drr, "F_{12}"] & & L_1 \otimes L_2 \arrow[dll, "F_{21}"'] \arrow[d, "F_{22}"] \\
E_1 \otimes E_2 \arrow[dr, "{\mathrm{id}_1 \otimes \psi_2,\ \psi_1 \otimes \mathrm{id}_2}"'] & & M_1 \otimes M_2 \arrow[dl, "{\varpi_1 \otimes \mathrm{id}_2,\ \mathrm{id}_1 \otimes \varpi_2}"] \\
 & (E_1 \otimes M_2) \times (M_1 \otimes E_2) = \Pi &
\end{tikzcd}

\note{le $\oplus$ de $\Sigma$ est écrit sur un signe biffé ; le $E_2$ du sommet $E_1 \otimes E_2$ de gauche, en bas, est surchargé.}

\marginal{\uncertain{car} \ill{} \uncertain{sommes} \ill{}}

Si on veut que $F$ induise \struck{l'identité} \add{un iso d'extensions, donc} $L_1 \otimes L_2 \to L_1 \otimes L_2$, il faut que $\boxed{F_{22} = 0}$, \emph{\uncertain{et on veut}} \add{de plus} que $\Sigma$ induise l'identité sur $L_1 \otimes L_2$, il faut que $\boxed{F_{21} = \alpha_1 \otimes \alpha_2}$, \uncertain{et pour qu'il} induise l'identité sur $M_1 \otimes M_2$, il faut que $\boxed{F_{12} = \varphi_1 \otimes \psi_2}$ : \struck{\ill{}} \uncertain{Dernière} condition \uncertain{sur la} \uncertain{vérification}, \uncertain{dans} \struck{\ill{}} des 4 cas \uncertain{voulus} \ill{} \uncertain{déjà} \uncertain{automatiques} (\uncertain{à droite} \uncertain{composés} \uncertain{nuls}) :

\note{la dernière égalité encadrée se lit $F_{12} = \varphi_1 \otimes \psi_2$ ; on attendrait $\psi_1 \otimes \psi_2$.}

\begin{tikzcd}[column sep=small, row sep=small]
 & \Sigma \arrow[dl] \arrow[dr] & \\
E_1 \otimes E_2 \arrow[drr, "{0,\ F_{12} = \psi_1 \otimes \psi_2}"'] & & L_1 \otimes L_2 \arrow[d, "F_{22} = 0"] \\
 & & M_1 \otimes M_2
\end{tikzcd}

et \struck{$E_1 \otimes E_2$}

\begin{tikzcd}[column sep=small, row sep=small]
L_1 \otimes L_2 \arrow[d, "F_{21}"'] \arrow[r, "F_{22} = 0"] & M_1 \otimes M_2 \arrow[dl] \\
E_1 \otimes E_2 \arrow[r] & \Pi
\end{tikzcd}

\note{disposition des deux derniers diagrammes reconstituée : l'auteur esquisse les composés qui doivent s'annuler ; le sens précis de certaines flèches est douteux.}

\page{167}
\begin{tikzcd}[column sep=small, row sep=small]
 & L_1 \otimes L_2 \arrow[dl] \arrow[dr] & \\
E_1 \otimes L_2 \arrow[dr] & \text{cocart.} & L_1 \otimes E_2 \arrow[dl] \\
 & E_1 \otimes_0 E_2 &
\end{tikzcd}

\begin{tikzcd}[column sep=small, row sep=normal]
0 \arrow[r] & E_1 \otimes_0 E_2 \arrow[r, "\mu"] \arrow[d, "{\pi_1 \otimes \mathrm{id},\ \mathrm{id} \otimes \pi_2 \ (\text{cocart.})}"'] & E_1 \otimes E_2 \arrow[r] \arrow[d] & M_1 \otimes M_2 \arrow[r] \arrow[d, no head, "\|" description] & 0 \\
0 \arrow[r] & L_1 \otimes L_2 \arrow[r] \arrow[d, no head, "\|" description] & E_1 \vee E_2 \arrow[r] \arrow[d, "?"] & M_1 \otimes M_2 \arrow[r] \arrow[d, no head, "\|" description] & 0 \\
0 \arrow[r] & L_1 \otimes L_2 \arrow[r] \arrow[d, no head, "\|" description] & E_1 \wedge E_2 \arrow[r] \arrow[d, "\text{cart.}"] & M_1 \otimes M_2 \arrow[r] \arrow[d, "{\varpi_1 \otimes \mathrm{id},\ \mathrm{id} \otimes \varpi_2}"] & 0 \\
0 \arrow[r] & L_1 \otimes L_2 \arrow[r] & E_1 \otimes E_2 \arrow[r] & E_1 \otimes^{\circ} E_2 \arrow[r] & 0
\end{tikzcd}

\note{un trait horizontal barre la page entre les deux premières lignes du diagramme, un « $E$ » biffé à côté de la flèche « ? » ; dans la dernière ligne, $E_1$ est écrit sur un $M_1$ biffé.}

\begin{tikzcd}[column sep=small, row sep=small]
 & E_1 \otimes^{\circ} E_2 \arrow[dl] \arrow[dr] & \\
E_1 \otimes M_2 \arrow[dr] & \text{cart.} & M_1 \otimes E_2 \arrow[dl] \\
 & M_1 \otimes M_2 &
\end{tikzcd}

\page{169}
\begin{tikzcd}[column sep=normal, row sep=small]
0 \arrow[r] & L_1 \arrow[r, "\alpha_1"] & E_1 \arrow[r, "\psi_1"] \arrow[l, bend right=40, dashed, "\pi_1"'] & M_1 \arrow[r] \arrow[l, bend right=40, dashed, "\varpi_1"'] & 0 \\
0 \arrow[r] & L_2 \arrow[r, "\alpha_2"] & E_2 \arrow[r, "\psi_2"] \arrow[l, bend left=40, dashed, "\pi_2"] & M_2 \arrow[r] \arrow[l, bend left=40, dashed, "\varpi_2"] & 0
\end{tikzcd}

$\pi_i$, $\varpi_i$ \uncertain{univoquement} \uncertain{déterminés} \uncertain{par} \uncertain{des} $\chi$-trivialisations

\[
(\pi_1, \varpi_1)\ (\pi_2, \varpi_2)
\left\{
\begin{array}{l}
\alpha_i \pi_i + \varpi_i \psi_i = \chi\, \mathrm{id}_{E_i} \\
\text{\struck{$\pi_1 \varpi_2 = \chi\, \mathrm{id}_{E_2}$}} \\
\pi_i \alpha_i = \chi\, \mathrm{id}_{L_i} \qquad \text{\struck{\ill{}}} = 0 \\
\psi_i \varpi_i = \chi\, \mathrm{id}_{M_i}
\end{array}
\right.
\]
\note{la dernière relation est écrite sur une formule biffée illisible.}

On définit de façon \uncertain{évidente} $E_1 \vee E_2$ \uncertain{et} $\pi$-trivialisations et $E_1 \wedge E_2$, extensions de $M_1 \otimes M_2$ par $L_1 \otimes L_2$, \uncertain{et} un \uncertain{isomorphisme} \uncertain{canonique} des extensions de $E_1 \vee E_2$ sur $E_1 \wedge E_2$
\[
\left[\;
\begin{array}{c}
E_1 \otimes L_2 \searrow \\
L_1 \otimes L_2 \qquad E_1 \otimes_0 E_2 \quad \text{cocartésien} \\
L_1 \otimes E_2 \nearrow
\end{array}
\;\right]
\]

\begin{tikzcd}[column sep=small, row sep=small]
0 \arrow[r] & E_1 \otimes_0 E_2 \arrow[r] \arrow[d, "{\pi_1 \otimes \mathrm{id},\ \mathrm{id} \otimes \pi_2}"'] & E_1 \otimes E_2 \arrow[r] \arrow[d] & M_1 \otimes M_2 \arrow[d, no head, "\|" description] \\
0 \arrow[r] & L_1 \otimes L_2 \arrow[r] & E_1 \vee E_2 \arrow[r] & M_1 \otimes M_2
\end{tikzcd}

\begin{tikzcd}[column sep=small, row sep=small]
0 \arrow[r] & L_1 \otimes L_2 \arrow[r] \arrow[d, no head, "\|" description] & E_1 \wedge E_2 \arrow[r] \arrow[d] & M_1 \otimes M_2 \arrow[d] \\
0 \arrow[r] & L_1 \otimes L_2 \arrow[r] & E_1 \otimes E_2 \arrow[r] & E_1 \otimes^{\circ} E_2
\end{tikzcd}

\note{dans la première ligne du second diagramme, le signe entre $E_1$ et $E_2$ est surchargé ($\wedge$ sur $\otimes$ ?).}

\section{Le cas des algèbres quadratiques}

\page{171}
\[
u = u_1 \star u_2 - b_1 b_2 = 2 u_1 u_2 + b_2 u_1 + b_1 u_2 \in \mathcal{A}_1 \otimes \mathcal{A}_2
\]
\note{un « ! » est écrit au-dessus du $\star$.}

Solution de l'équation
\[
u^2 + bu + c = 0 \quad \text{dans } \mathcal{A}
\]
(\uncertain{conjugués} \ill{} car. $2$ si $b_1 = b_2 = 0$)
\[
b = b_1 b_2, \qquad c = b_1^2 c_2 + b_2^2 c_1 - 4 c_1 c_2
\]
\[
\delta = b^2 - 4c = \underbrace{\delta_1}_{(b_1^2 - 4c_1)} \underbrace{\delta_2}_{(b_2^2 - 4c_2)}
\]

\begin{tikzcd}[column sep=normal, row sep=small]
\mathcal{A}_1 \otimes \mathcal{A}_2 \arrow[r, "\text{com. par div. à } \mathcal{A}"] \arrow[rr, bend left=30, "\sigma"] & \mathcal{A}_1 \wedge \mathcal{A}_2 = \mathcal{A} \arrow[r] & \mathcal{A}_1 \otimes \mathcal{A}_2 \\
E_1 \otimes E_2 & E = E_1 \wedge E_2 \supset L_1 \otimes L_2 \arrow[l] \arrow[d] & E_1 \otimes E_2 \arrow[l, "\text{com. diff. à } E"'] \arrow[d] \arrow[ll, bend left=30, "\sigma"'] \\
 & 0 & 0
\end{tikzcd}

\note{diagramme recomposé : deux lignes, l'algèbre en haut et le module en bas, chacune surmontée d'un arc marqué d'un symbole de type $\sigma$ (au-dessus de $\mathcal{A}$, « $\underline{O}$ » souligné) ; des flèches verticales descendent vers $0$. Les étiquettes « com. par div. à $\mathcal{A}$ » et « com. diff. à $E$ » sont d'une lecture douteuse.}

\[
f \in \mathcal{A}_1 \otimes \mathcal{A}_2 \longmapsto f^{\sigma} \;\text{\struck{$(\sigma_1 \otimes \sigma_2 + \mathrm{id})(f)$}}
\]
\[
= \bigl(\underbrace{\sigma_1 \otimes \sigma_2}_{\sigma} + \mathrm{id}_{\mathcal{A}_1 \otimes \mathcal{A}_2}\bigr)(f) = f + \bar f
\]
\[
x \in E_1 \otimes E_2 \qquad x^{\sigma} = x + \bar x
\]

\begin{tikzcd}[column sep=normal, row sep=small]
\mathcal{A}_1 \otimes_0 \mathcal{A}_2 \arrow[r, "\mathrm{Tr}"] \arrow[d] & \underline{O} \arrow[r, "\mathrm{id}"] \arrow[d] & \underline{O} \arrow[d] \\
\mathcal{A}_1 \otimes \mathcal{A}_2 \arrow[r] \arrow[d] & \mathcal{A} \arrow[r] \arrow[d] & \mathcal{A}_1 \otimes \mathcal{A}_2 \arrow[d] \\
L_1^{\vee} \otimes L_2^{\vee} \arrow[r, "\mathrm{id}"] \arrow[d] & L_1^{\vee} \otimes L_2^{\vee} \arrow[r, "2"] \arrow[d] & L_1^{\vee} \otimes L_2^{\vee} \arrow[d] \\
0 & 0 & 0
\end{tikzcd}

\note{dans la colonne de droite, $\underline{O}$ est suivi de $\mathcal{A}_1 \otimes_0 \mathcal{A}_2$ avant $\mathcal{A}_1 \otimes \mathcal{A}_2$ ; la disposition est simplifiée ici.}

\struck{$(\mathcal{A}_1 \otimes \mathcal{A}_2)^{\sigma}$ \ill{} \uncertain{définit} \ill{} car. $2$}

\[
0 \to L_1 \to E_1 \to L_1' \to 0, \qquad 0 \to L_2 \to E_2 \to L_2' \to 0
\]
$E_1 \otimes E_2$ \qquad $E_1 \otimes E_2 \to L_1 \otimes L_2$ \qquad $L_1 \otimes E_2$

\page{173}
\[
0 \to L \xrightarrow{\;i\;} E \xrightarrow{\;\varphi\;} L' \to 0
\]
\[
E \xrightarrow{\;\pi\;} L \quad \text{tel que } \pi(x) = \chi x \text{ si } x \in L
\]
\[
L' \xrightarrow{\;\varpi\;} E \quad \text{tel que } \varphi(\varpi(\lambda)) = \chi \lambda \text{ si } \lambda \in L'
\]
$E \xrightarrow{f} E$ tel que
\begin{enumerate}
\item[1°)] $f$ nul sur $L$
\item[2°)] \struck{$f(x) - \chi x \in L$} $\chi x - f(x) \in L$ \quad \struck{$f - \chi\,\mathrm{id}$} \quad $\pi(x)$
\end{enumerate}
\uncertain{donc} $E \xrightarrow{\varphi} L' \xrightarrow{\varpi} E$ \ill{}

\[
\boxed{\pi = \chi\,\mathrm{id} - f}
\]

\struck{$\pi$ \uncertain{identité} dans $L$ $\Longleftrightarrow$}

$\varpi$ \uncertain{stabilise} $L$ \quad $\longleftrightarrow$ \uncertain{$f$ nul sur} $L \hookrightarrow E/L$ \ill{}

$\pi$ : \uncertain{id.} dans $L$ $\longleftrightarrow$ $f$ \uncertain{induit} \struck{\ill{}} $\chi\,\mathrm{id}_{L'}$ sur $E/L$

$\pi$ \uncertain{induit} $\chi\,\mathrm{id}_L$ sur $L$ $\longleftarrow$ $f$ nul sur $L'$

\note{ces trois lignes d'équivalences sont reliées par une accolade et un trait ; leur lecture est douteuse.}

\[
0 \to L^{\vee} \hookrightarrow \mathcal{A} \to L^{\vee} \to 0
\]
\[
0 \to L \to E_n \to L' \to 0
\]
\struck{$0 \to L \otimes M \to E \otimes M \to L' \otimes M \to 0$}
\[
0 \to L_1 \to E_1 \to L_1' \to 0 \qquad \text{\struck{$E_1 \otimes E_2$}} \qquad 0 \to L_2 \to E_2 \to L_2' \to 0
\]
\[
0 \to L_1 \to E_1 \to L_1' \to 0
\]
\[
0 \to L_2 \to E_2 \to L_2' \to 0 \qquad E_1 \simeq \mathcal{A}_1 \otimes \Delta_1
\]
\[
E_1 \wedge E_2 \qquad
\overset{\Delta_1^{-1}}{E_1^{\vee}} \wedge \overset{\Delta_2^{-1}}{E_2^{\vee}}
\simeq \bigl(\underbrace{E_1 \wedge E_2}_{\Delta_1 \Delta_2}\bigr) \otimes \Delta_1^{-1} \otimes \Delta_2^{-1}
\]
\[
\simeq (E_1 \wedge E_2)^{\vee}
\]

\[
\pi(x) = 2x - F(x)
\]
\[
\pi(x) - x = x - F(x)
\]
\note{le $F$ de la première ligne est écrit sur un $f$.}

\page{175}
\[
\begin{array}{llcll}
\mathcal{A}_i & & & \mathcal{A} & u \qquad 1 \\[4pt]
E_i & x_i,\ e_i & & E & x_1 \star x_2 = x,\quad e = e_1 \otimes e_2 \\
 & \cap \quad \cap & & & \cap \\
 & P_i \quad L_i & & & P_3
\end{array}
\]
\note{l'indice « $3$ » de $P_3$ est douteux ; le $E$ de la colonne de droite est écrit sur une lettre biffée.}

\[
\begin{array}{ll}
\langle u_i, x_i \rangle = \langle 1_i, e_i \rangle = 0 & \qquad \langle u, e \rangle = \langle 1, x \rangle = 1 \\
\langle u_i, e_i \rangle = \langle 1_i, x_i \rangle = 1 & \qquad \langle u, x \rangle = \langle 1, e \rangle = 0
\end{array}
\]

$u_1 \in \mathcal{A}_1$, \quad $u_2 \in \mathcal{A}_2$ \qquad \struck{$\mathcal{A}_1 \otimes \mathcal{A}_2$} \quad On veut exprimer

\struck{$\sigma(u_1 \star u_2) = 2 u_1 u_2 + b_2 u_1 +$}

$u_1 \star u_2$ comme $\lambda 1_{\mathcal{A}} + \mu u$, \ill{} \uncertain{calculs}
\[
\lambda = \langle \underbrace{u_1 \star u_2}_{u}, x \rangle, \qquad \mu = \langle u_1 \star u_2, e \rangle
\]
\struck{$\mathcal{A}_1 \otimes \mathcal{A}_2 \to \mathcal{A}$}
\[
\langle u_1 \star u_2, x_1 \otimes x_2 \rangle
= \langle \underbrace{u_1 \otimes u_2 + \bar u_1 \otimes \bar u_2}_{(u_1 \otimes u_2)^{\sigma}}, x_1 \otimes x_2 \rangle
\]
\[
= \langle 2 u_1 \otimes u_2 + b_1 u_2 + u_1 b_2 + b_1 b_2,\ x_1 \otimes x_2 \rangle
\]
\[
= b_1 b_2
\]
\note{un « ! » sous le terme $b_1 u_2$ ; un début de ligne « $= 2\beta_1 + \beta$ » est biffé, ainsi qu'un terme illisible devant $2 u_1 \otimes u_2$.}

\[
\mu = \langle u_1 \star u_2, e \rangle
\]
\struck{$= \langle \varphi_1(u_1).e_1 \rangle \langle \varphi_2(u_2), e_2 \rangle$}

\struck{$= \langle 2 u_1 \otimes u_2 + b_1 u_2 + u_1 b_2 + b_1 b_2,\ e_1 \otimes e_2 \rangle$}
\[
= 1 \qquad e \in T\,\mathcal{L} \subset T\!E
\]
\note{lecture de « $e \in T\mathcal{L} \subset TE$ » très douteuse.}

donc
\[
u_1 \star u_2 = b_1 b_2 + \text{\struck{\ill{}}}\ u
\]
i.e.
\[
\boxed{u_1 \star u_2 - b_1 b_2 = u}
\]
\[
u_1 \star u_2 - \mathrm{Tr}_1(u_1)\,\mathrm{Tr}_2(u_2) = 2 \,\cdot
\]
L'image de $u$ dans $\mathcal{A}_1 \otimes \mathcal{A}_2$ est donc
\[
\boxed{2 u_1 \otimes u_2 + b_2 u_1 + b_1 u_2}
\]

\page{177}
Compatibilités pour

\begin{tikzcd}[column sep=normal, row sep=small]
\mathcal{A}_1, \mathcal{A}_2 \arrow[r, "\rightrightarrows" description] & \mathcal{A}_1 \otimes_0 \mathcal{A}_2 \arrow[r, "\mathrm{Tr}"] \arrow[dr, hook] & \underline{O} \arrow[r] & \mathcal{A} \\
 & & \mathcal{A}_1 \otimes \mathcal{A}_2 \arrow[ur] &
\end{tikzcd}

\begin{tikzcd}[column sep=normal, row sep=small]
{E_1 \otimes L_2,\ L_1 \otimes E_2} \arrow[r, "\rightrightarrows" description] & E_1 \otimes_0 E_2 \arrow[r, "\pi"] \arrow[dr, hook] & L_1 \otimes L_2 \arrow[r] & E \\
 & & E_1 \otimes E_2 \arrow[ur] &
\end{tikzcd}

\note{une flèche biffée part de $L_1 \otimes L_2$ ; la flèche vers $E$ part de $L_1 \otimes L_2$ en oblique.}

Il y a à écrire quatre compatibilités
\begin{enumerate}
\item[1°)] si $u \in L_1 \otimes L_2$, $x \in E_1 \otimes E_2$, \uncertain{donc} $u \oplus x$ \uncertain{donne} … ;
$\underline{O} \oplus (\mathcal{A}_1 \otimes \mathcal{A}_2) \to \underline{O}$ qui donne … $=$ \uncertain{donc} \struck{\ill{}} sur $\mathcal{A}_1$.
\item[2°)] $f \in \mathcal{A}_1 \otimes_0 \mathcal{A}_2$ alors
\[
\langle \underbrace{f - f^{\sigma}}_{\mathrm{Tr}(f).1}, \text{\ill{}} \rangle = \mathrm{Tr}(f)\, \langle 1 \otimes \text{\ill{}}, x \rangle
\]
\item[3°)] $x \in E_1 \otimes E_2$
\[
\langle f, x + \bar x \rangle = \langle f, \pi(x) \rangle = \varphi(f) . \pi(x)
\]
\end{enumerate}

OK
\note{la page s'arrête sur « OK », entre deux traits ; le quatrième point annoncé n'est pas écrit ici.}

\page{179}
\[
\mathcal{A} \overset{\mathrm{def}}{=} (\mathcal{A}_1 \otimes \mathcal{A}_2) \oplus_{(\mathcal{A}_1 \otimes_0 \mathcal{A}_2)} \underline{O}
\]
pour

\begin{tikzcd}[column sep=normal, row sep=small]
 & & \mathcal{A}_1 \otimes \mathcal{A}_2 \\
\mathcal{A}_1 \otimes_0 \mathcal{A}_1 = \mathcal{A}_1 \otimes_{\underline{O}} \mathcal{A}_2 \arrow[urr, hook] \arrow[drr, "{\mathrm{Tr}_1, \mathrm{Tr}_2}"'] & & \\
 & & \underline{O}
\end{tikzcd}

\note{l'indice du second facteur de $\mathcal{A}_1 \otimes_0 \mathcal{A}_1$ est lu tel quel ; une flèche biffée et un signe biffé précèdent l'inclusion.}

\[
E \overset{\mathrm{def}}{=} (E_1 \otimes E_2) \oplus_{E_1 \otimes_0 E_2} L_1 \otimes L_2
\]
pour

\begin{tikzcd}[column sep=small, row sep=small]
E_1 \otimes_0 E_2 \arrow[r, hook] \arrow[d, no head, "\|" description] & E_1 \otimes E_2 \\
E_1 \otimes L_2 \oplus_{L_1 \otimes L_2} L_1 \otimes E_2 \arrow[r, "{(\pi_1 \otimes \mathrm{id},\ \mathrm{id} \otimes \pi_2)}"] & L_1 \otimes L_2
\end{tikzcd}

\uncertain{Accouplement}, pour \ill{} le \uncertain{produit} \uncertain{scalaire} (\uncertain{compatible} \ill{} \uncertain{essentiellement} \ill{}
\[
\sigma = \sigma_1 \otimes \sigma_2 \quad \text{dans } \mathcal{A}_1 \otimes \mathcal{A}_2 \text{ et } E_1 \otimes E_2,
\]
\ill{} \uncertain{sont} \uncertain{transposés} l'un de l'autre,
\[
f \longmapsto f^{\sigma} = f + \sigma f, \quad f \in \mathcal{A}_1 \otimes \mathcal{A}_2, \qquad x^{\sigma} = x + \sigma x, \quad x \in E_1 \otimes E_2
\]

\struck{$f, x$}
\[
\langle f, x^{\sigma} \rangle = \langle f^{\sigma}, x \rangle = \langle f, x^{\sigma} \rangle
\]
i.e.
\[
\langle f_1 \otimes f_2, \underbrace{x_1 \otimes x_2}_{?} \rangle
= \bigl\langle f_1 \otimes f_2 + (\mathrm{Tr}_1(f_1) - f_1)(\mathrm{Tr}_2(f_2) - f_2),\ x_1 \otimes x_2 \bigr\rangle
\]
\[
= \bigl\langle f_1 \otimes f_2,\ (\pi_1(x_1)T_1 - x_1) \otimes (\pi_2(x_2)T_2 - x_2)
\]
\[
\qquad + x_1 \otimes x_2 \bigr\rangle
\]
\note{la seconde égalité est d'une lecture douteuse (« $\pi_1(x_1)T_1$ » écrit sur une lettre biffée).}

\struck{i.e. $\langle f_1, x_1 \rangle \langle f_2, x_2 \rangle +$ \ill{} \quad $f_1 \otimes f$ c'est d'\uncertain{autres}}

Sur les \uncertain{deux} \ill{} \ill{} on aura
\[
\begin{array}{ll}
\langle \lambda, u \rangle = 0 & \text{si } \lambda \in \underline{O},\ u \in L_1 \otimes L_2 \\[2pt]
\langle f, u \rangle = \langle \varphi(f), u \rangle = \varphi(f).u & \text{si } f \in \mathcal{A}_1 \otimes \mathcal{A}_2,\ u \in L_1 \otimes L_2 \\[2pt]
\langle \lambda, x \rangle = \lambda\, \pi(x) & \text{si } x \in E_1 \otimes E_2,\ \lambda \in \underline{O}
\end{array}
\]
\note{la feuille s'arrête là ; l'accouplement se poursuit vraisemblablement au-delà du lot.}

\end{document}
